\documentclass[12pt,a4paper]{article}
\usepackage[UTF8]{ctex}
\usepackage{geometry}
\geometry{left=2.5cm,right=2.5cm,top=2.5cm,bottom=2.5cm}
\usepackage{graphicx}
\usepackage{amsmath,amsfonts,amssymb}
\usepackage{booktabs}
\usepackage{multirow}
\usepackage{array}
\usepackage{caption}
\usepackage{subcaption}
\usepackage{float}
\usepackage{xcolor}
\usepackage{hyperref}
\usepackage{natbib}
\usepackage{indentfirst}
\setlength{\parindent}{2em}

\title{\textbf{面向多尺度时序预测的时域门控混合模型}}
\author{付文浩$^{1}$，李荣丽$^{2,*}$\\
\small $^{1}$云南师范大学数学学院，昆明650500\\
\small $^{2}$云南师范大学数学学院，昆明650500}
\date{}

\begin{document}
\maketitle

\begin{abstract}
针对现有时间序列预测模型难以同时兼顾短期预测的高时效响应与长期预测的高精度需求这一问题，本文提出了一种时域门控混合模型（HybridModel）。该模型引入双维度时域门控融合单元，基于物理先验固定调节DLinear线性路径与TimeXer深度注意力路径的融合权重，实现短期惯性主导区与长期非线性演变区的结构化切换；同时设计尺度自适应损失函数，协调多时间尺度下的训练目标冲突。基于空气质量、电力负荷和气象等多领域数据的实证结果表明：在PM2.5预测任务中，HybridModel在三城市四步长全部12组实验中均优于TimeXer基线（提升2.8\%--8.9\%）；推理延迟仅1.81\,ms，满足实时业务需求。模型参数量较TimeXer仅增加1.9\%，在精度与效率间实现了有效权衡。

\textbf{关键词：}时域门控；混合模型；多尺度预测；时间序列预测；深度学习
\end{abstract}

\begin{center}
\textbf{\large Time-Scale-Aware Gating Hybrid Model for Multiscale Time Series Forecasting}

\vspace{0.5em}
\noindent\textbf{Abstract:} To address the dilemma between high-efficiency short-term responsiveness and high-accuracy long-term forecasting in time series prediction, this study proposes a HybridModel with dual-dimension temporal gating. The model structurally fuses a DLinear linear path and a TimeXer deep attention path via a time-dimension gate based on physical priors, while a variable-dimension gate enables multi-factor fusion. A scale-adaptive loss function is designed to harmonize divergent optimization objectives across temporal scales. Empirical validation on air quality, electricity, and weather datasets demonstrates that HybridModel outperforms TimeXer baseline in all 12 settings (improvement 2.8\%--8.9\%), with inference latency of only 1.81\,ms. The model increases parameters by merely 1.9\% over TimeXer, achieving an effective trade-off between accuracy and efficiency.

\noindent\textbf{Keywords:} temporal gating; hybrid model; multiscale prediction; time series forecasting; deep learning
\end{center}

\section{引言}

时间序列预测在环境监测、电力调度、金融风控等业务场景中具有广泛应用。以PM2.5浓度预测为例，突发性污染过程的应急决策需要模型具备毫秒级响应的``短期迅捷''，而中长期环境治理规划则依赖于对污染物演化趋势的``长期精准''。然而，现有深度学习预测模型往往陷入``精度''与``效率''的二元悖论：以TimeXer\cite{wang2024timexer}、iTransformer\cite{liu2023itransformer}为代表的高精度模型因注意力机制计算开销大，推理延迟高；以DLinear\cite{zeng2023transformers}、TSMixer\cite{chen2023tsmixer}为代表的轻量级模型虽效率卓越，但模型容量简单，长期预测中误差累积显著。这种``高精度模型失之于迟，高效率模型失之于拙''的矛盾，迫使现有业务系统往往需维护多套模型，增加了部署复杂性与运维成本。

现有混合策略多聚焦于输入表征层面的分解与融合（如Autoformer\cite{wu2021autoformer}的序列分解、TimeMixer\cite{wang2024timemixer}的多尺度混合），鲜有研究从输出层面进行跨架构的结构化融合。更为关键的是，多数混合模型采用数据驱动的自适应门控，其融合权重依赖梯度学习，但在主路径（DLinear/TimeXer）梯度主导下，门控参数往往难以有效学习，最终退化为固定加权平均。

针对上述挑战，本文提出了一种面向多尺度时序预测的时域门控混合模型（HybridModel）。该模型的核心创新在于：（1）\textbf{物理驱动的结构化融合}——基于时间序列的周期物理先验（如日周期边界），在输出层固定划分短期/长期预测路径，而非依赖数据学习；（2）\textbf{双维度时域门控机制}——时间维度门控依据物理规律切换DLinear线性路径与TimeXer注意力路径，变量维度门控实现多因子的差异化融合；（3）\textbf{尺度自适应损失函数}——通过门控加权协调短期精度与长期鲁棒性的训练目标冲突。在PM2.5空气质量预测任务上的实证表明，HybridModel以仅1.9\%的参数量增量和1.81\,ms的推理延迟，在三城市四步长全部12组实验中均优于TimeXer基线，验证了该融合策略在精度-效率权衡上的实用价值。

\section{相关工作}

\subsection{PM2.5预测方法演进}

PM2.5浓度预测经历了从统计模型到深度学习的发展过程\cite{lim2021time}。早期研究主要依赖ARIMA、指数平滑等统计方法，这些方法在平稳序列上表现良好，但难以处理大气污染数据中的非线性和非平稳特征\cite{box2015time}。随着机器学习的发展，随机森林、支持向量回归和梯度提升树等方法被引入空气质量预测，显著提升了非线性拟合能力。然而，这些方法对特征工程的依赖较强，且对长期时间依赖的建模能力有限\cite{mendez2023machine}。

近年来，深度学习已成为PM2.5预测的主流范式。Qi等\cite{qi2019hybrid}将图卷积网络与LSTM结合，实现了PM2.5的时空预测；Zhu和Xie\cite{zhu2023investigation}采用1D-CNN-BiLSTM并行结构进行多站点PM2.5预报；Xia等\cite{xia2024multimodal}则融合遥感图像与时间序列数据，构建了多模态深度学习框架。在Transformer架构方面，Zhang和Zhang\cite{zhang2023modeling}提出了稀疏注意力Transformer网络用于PM2.5预测；Tong等\cite{tong2024robust}设计了基于Transformer的时空预测模型；Feng等\cite{feng2023spatiotemporal}将Informer架构应用于空气质量预报；Zeng等\cite{zeng2023ceemdan}采用CEEMDAN分解与深度Transformer相结合进行长期PM2.5浓度预测。这些研究表明，深度学习在捕捉大气污染复杂时空模式方面具有显著优势\cite{su2025systematic}。

尽管上述研究在单一模型架构上取得了进展，但鲜有工作系统性地探讨不同时间尺度下模型架构的适配性问题——即短期预测中轻量级模型的效率优势与长期预测中深度模型的精度优势如何有机融合。

\subsection{通用时间序列预测模型进展}

在时间序列预测领域，Transformer架构近年来取得了突破性进展。Wu等\cite{wu2021autoformer}提出的Autoformer通过自相关机制和序列分解实现了长期预测；Zhou等\cite{zhou2022fedformer}在FEDformer中引入频域增强分解；Woo等\cite{woo2022etsformer}设计了指数平滑Transformer。针对多变量预测，Liu等\cite{liu2023itransformer}提出的iTransformer将注意力机制应用于变量维度而非时间维度；Wang等\cite{wang2024timexer}的TimeXer通过patch-wise嵌入和交叉注意力机制实现了内源变量与外源变量的联合建模。

与此同时，研究者也开始探索Transformer之外的轻量级架构。Zeng等\cite{zeng2023transformers}的DLinear通过简单的趋势-季节分解和线性映射，在多个基准上超越了复杂Transformer模型；Wu等\cite{wu2023timesnet}的TimesNet将时间序列转换为2D空间进行多周期建模；Nie等\cite{nie2022time}的PatchTST通过通道独立的patch设计扩展了注意力机制的长程覆盖能力；Chen等\cite{chen2023tsmixer}的TSMixer采用全MLP架构实现高效的多尺度时序预测。

然而，现有模型大多采用单一架构范式，要么追求精度（如TimeXer）而牺牲效率，要么追求效率（如DLinear）而损失长期预测能力。如何在同一框架内协调这两种范式的优势，仍是时序预测领域的重要开放问题。

\subsection{混合模型与多尺度融合策略}

混合模型通过整合多种架构的优势以提升预测性能。在空气质量预测领域，混合策略主要体现为``分解-预测''框架和``多模型集成''两类。Xie等\cite{xie2024daily}利用变分模态分解与深度学习进行多站点时空融合；Ding和Sun\cite{ding2023prediction}设计了加权RF-LSTM混合模型。在通用时序预测领域，Wang等\cite{wang2024timemixer}的TimeMixer引入多尺度可分解混合机制，在不同时间分辨率上进行信息交互。

上述混合策略多聚焦于输入表征层面的分解与融合，而本文提出的HybridModel则从输出层面进行跨模型融合——直接结合DLinear的线性外推与TimeXer的注意力表征，并基于大气物理规律设计固定的时间尺度划分策略。这种``物理驱动+结构化融合''的设计思路，在空气质量数据强日周期特性的场景下展现出独特优势。
\section{方法论}

\subsection{大气污染的多尺度物理分解}

大气PM2.5浓度的演变是一个多物理过程耦合的系统。根据大气边界层理论\cite{stull1988introduction}，我们将PM2.5的时间序列分解为两个正交的物理分量：

\textbf{高频湍流分量}（High-Frequency Turbulence Component）：主要由物理惯性（如污染气团平流输送）主导，变化近似线性。模型在短期预测时倾向于启用DLinear路径，利用其高效的线性外推能力快速捕捉瞬时波动，满足实时预警的低延迟需求。

\textbf{低频趋势分量}（Low-Frequency Trend Component）：受化学演变（如光化学反应生成二次颗粒物\cite{huang2014high}）及大尺度天气过程影响，呈现复杂的非线性特征。此时，模型自动增强TimeXer路径的权重，借助其深度注意力机制捕捉长期依赖与非线性转化过程，以提升预测精度。

\subsection{双维度时域门控融合单元}

传统门控机制多作用于单模型内部（如AdaptiveGate\_TimeXer在注意力层内部剪枝），其门控输出为注意力模块与前馈网络之间的融合权重。本文提出一种\textbf{跨模型门控策略}，直接在DLinear的线性外推输出与TimeXer的深度表征输出之间进行基于物理先验的融合，使模型能够根据预测步长的时间尺度特性选择最优计算路径。

\subsubsection{时间维度门控（Time Gate）}

设DLinear输出为 $\mathbf{y}_D \in \mathbb{R}^{B \times L \times D}$，TimeXer输出为 $\mathbf{y}_T \in \mathbb{R}^{B \times L \times D}$，其中 $B$ 为批次大小，$L$ 为预测步长，$D$ 为变量数。

定义时间维度门控为基于物理先验的固定参数经Sigmoid变换：
\begin{equation}
\mathbf{g}_t = \sigma(\mathbf{w}_t) \in \mathbb{R}^{L}
\end{equation}
其中 $\mathbf{w}_t \in \mathbb{R}^{L}$ 为可学习参数，$\sigma(\cdot)$ 为Sigmoid函数。

初始化策略基于物理先验：前 $s_0=12$ 步对应大气边界层的一个典型日变化周期，此期间物理惯性主导，故令 $\mathbf{w}_t[:s_0] = 2.0 \Rightarrow \mathbf{g}_t \approx 0.88$（偏向DLinear）；超过日周期后化学演变引入非线性，令 $\mathbf{w}_t[s_0:] = -2.0 \Rightarrow \mathbf{g}_t \approx 0.12$（偏向TimeXer）。该硬边界初始化为门控提供物理引导，在训练过程中保持固定以维持稳定的时间尺度划分。

\subsubsection{变量维度门控（Variable Gate）}

时间门控仅实现了``全局时间尺度的路径切换''，但未考虑不同污染因子在同一时刻的差异化重要性。为此，设计变量维度门控：
\begin{equation}
\mathbf{g}_v = \sigma\bigl(\text{MLP}([\mathbf{y}_D \| \mathbf{y}_T \| \mathbf{g}_t^{\text{expand}}])\bigr) \in \mathbb{R}^{B \times L \times D}
\end{equation}
其中 $[\cdot\|\cdot]$ 表示通道拼接，$\mathbf{g}_t^{\text{expand}}$ 将时间门控广播至 $(B, L, 1)$ 后与输出拼接。MLP采用两层线性映射：$\text{MLP}(\mathbf{x}) = \mathbf{W}_2 \cdot \text{GELU}(\mathbf{W}_1 \mathbf{x} + \mathbf{b}_1) + \mathbf{b}_2$，隐层维度为 $D/4$。

\subsubsection{融合输出}
\begin{equation}
\mathbf{y}_{\text{fused}} = \mathbf{g}_v \odot \mathbf{y}_D + (1 - \mathbf{g}_v) \odot \mathbf{y}_T
\end{equation}
其中 $\odot$ 为逐元素乘。最终输出经Dropout正则化后传递至下游。

\subsection{HybridModel整体架构}

HybridModel由三个核心组件构成（图~\ref{fig:architecture}示意）：

\textbf{DLinearComponent}：对输入序列进行``趋势--季节''分解后分别通过单层线性映射外推，输出 $\mathbf{y}_D$。该组件计算复杂度为 $\mathcal{O}(L \cdot D)$，适用于短期惯性主导区。

\textbf{TimeXerComponent}：采用patch-wise embedding将序列切分为 patches，通过编码器（自注意力+交叉注意力）捕捉变量间耦合关系，输出 $\mathbf{y}_T$。该组件计算复杂度为 $\mathcal{O}(L^2 \cdot D)$，适用于长期非线性演变区。

\textbf{AdaptiveGateUnit（时域门控融合单元）}：接收 $\mathbf{y}_D$ 和 $\mathbf{y}_T$，通过时间门控和变量门控生成融合输出 $\mathbf{y}_{\text{fused}}$。门控单元参数量仅122个，可忽略不计；训练过程中保存门控值用于可视化分析。

从计算效率角度，HybridModel单次前向推理需依次执行DLinear和TimeXer两条路径及门控融合，实测推理延迟为$\approx$1.81\,ms（batch\_size=32, pred\_len=96），较TimeXer单路径（$\approx$1.32\,ms）增加约38\%，较iTransformer（$\approx$0.82\,ms）增加约121\%（图~\ref{fig:latency}）。尽管延迟有所增加，但在环境监测业务中，毫秒级差异对实时预警不构成瓶颈；相反，HybridModel以接近TimeXer的延迟（差异$<$1\,ms）获得了DLinear路径的短期预测增益，在精度-效率权衡上具有实用价值。未来可通过将两条主路径并行化（GPU多流执行）进一步压缩延迟至$\max(T_D, T_T)$水平。

\subsection{尺度自适应损失函数}

普通MSE损失对所有预测步长和变量赋予同等权重，无法平衡长期趋势预测（需鲁棒性）和短期波动预警（需敏感度）之间的矛盾。本文设计尺度自适应损失函数（Scale-Adaptive Loss, SAL）：
\begin{equation}
\mathcal{L}_{\text{SAL}} = \frac{1}{BLD} \sum_{b,l,d} \left[ g_v^{(b,l,d)} \cdot (y_D^{(b,l,d)} - \hat{y}^{(b,l,d)})^2 + (1 - g_v^{(b,l,d)}) \cdot (y_T^{(b,l,d)} - \hat{y}^{(b,l,d)})^2 \right]
\end{equation}
该损失使DLinear路径的误差在高门控区域（短期）被放大，TimeXer路径的误差在低门控区域（长期）被放大，从而驱动两个分支分别专精于其优势区间。若门控值缺失（如eval模式），则退化为加权平均：$\mathcal{L} = \alpha \cdot \mathcal{L}_{\text{MSE}}(y_D, \hat{y}) + (1-\alpha) \cdot \mathcal{L}_{\text{MSE}}(y_T, \hat{y})$，其中 $\alpha=0.7$。

\begin{figure}[H]
\centering
\includegraphics[width=0.85\textwidth]{figures/inference_latency.png}
\caption{各模型推理延迟对比（batch\_size=32, seq\_len=96, pred\_len=96, RTX 4090）。误差线为标准差（100次运行）。}
\label{fig:latency}
\end{figure}

\section{数据说明及实验设置}

\subsection{实验数据}

\subsubsection{昆明市数据（清洁基准）}

数据来源于RESSET空气质量监测平台。研究选取昆明市2019-02-18 17:00至2025-01-23 00:00的空气质量数据，共计46,573小时。使用6个指标（PM10、SO2、CO、NO2、O3、PM2.5）做多元时间序列预测，预测目标为PM2.5。PM2.5浓度均值为22.3 $\mu$g/m$^3$，标准差12.8，主要分布于Excellent（0--35）和Good（35--75）等级。数据按7:2:1划分为训练集（32,601）、验证集（4,657）和测试集（9,315）。

\subsubsection{北京市数据（重污染验证）}

数据来源于Open-Meteo Air Quality API（ECMWF CAMS再分析数据），时段为2022-08至2023-12，共计约12,000小时。PM2.5浓度均值为71.2 $\mu$g/m$^3$，标准差45.6，覆盖Excellent至Heavy（$>$150）全部五个污染等级。该数据集用于验证模型在重污染、高波动场景下的泛化能力。划分比例同昆明。

\subsubsection{上海市数据（中等污染过渡）}

数据来源同北京，时段为2022-08至2023-12，约12,000小时。PM2.5浓度均值为45.8 $\mu$g/m$^3$，标准差28.3，主要分布于Excellent、Good和Light（75--115）等级。该数据集作为污染梯度上的过渡点，检验模型在不同污染强度间的平滑适配性。划分比例同昆明。

\subsection{变量选择}

基于皮尔逊相关系数分析，PM10与PM2.5存在强相关（$r=0.90$），能直接反映主要污染来源；CO和NO2的中等相关性（$r=0.63/0.54$）捕捉燃烧过程和机动车排放的二次转化；SO2的弱相关（$r=0.40$）表征工业源贡献；O3虽与PM2.5直接相关性弱，但其参与的光化学反应生成二次颗粒物\cite{huang2014high}，间接影响PM2.5生成速率。上述6个变量共同构建了覆盖扬尘、工业、机动车、光化学四大污染机制的完整输入体系。

\subsection{参数设置与对比模型}

实验基于TSLib（Time-Series-Library）框架实现。输入序列长度 $L_{\text{in}}=96$，预测步长 $L_{\text{out}} \in \{96, 192, 336, 720\}$（对应4天、8天、14天、30天）。所有数据经StandardScaler归一化处理。

对比模型包括：（1）TimeXer\cite{wang2024timexer}：patch-wise Transformer基线；（2）iTransformer\cite{liu2023itransformer}：倒置Transformer；（3）VCGU\_iTtransformer：轻量化变体；（4）AdaptiveGate\_TimeXer：带门控的TimeXer改进版；（5）HybridModel（本文）：时间尺度感知门控混合模型。各模型核心参数见表~\ref{tab:params}。

\begin{table}[H]
\centering
\caption{各模型核心参数与计算复杂度}
\label{tab:params}
\begin{tabular}{lccccc}
\toprule
模型 & d\_model & e\_layers & 参数量 & 模型大小 & 时间复杂度 \\
\midrule
DLinear & -- & -- & 1.9$\times10^4$ & 0.07MB & $\mathcal{O}(L \cdot D)$ \\
TimeXer & 256 & 1 & 9.9$\times10^5$ & 3.79MB & $\mathcal{O}(L^2 \cdot D)$ \\
iTransformer & 512 & 3 & 6.5$\times10^6$ & 24.8MB & $\mathcal{O}(D^2 \cdot L)$ \\
VCGU\_iTtransformer & 128 & 2 & 8.2$\times10^5$ & 3.13MB & $\mathcal{O}(L \cdot D)$ \\
AdaptiveGate\_TimeXer & 256 & 1 & 1.0$\times10^6$ & 3.82MB & $\mathcal{O}(L^2 \cdot D)$ \\
\textbf{HybridModel} & 256 & 1 & \textbf{1.0$\times10^6$} & \textbf{3.86MB} & $\mathcal{O}(L^2 \cdot D)$ \\
\bottomrule
\end{tabular}
\end{table}

\textbf{注：}参数量基于pred\_len=96、enc\_in=6的配置统计；HybridModel较TimeXer仅增加约1.9\%的参数量（18,746个），时间复杂度主要由TimeXer路径的交叉注意力机制主导。DLinear参数量极小但长期预测精度受限。HybridModel以接近TimeXer的模型规模，实现了双路径的物理结构化融合。

\section{实验结果与分析}

\subsection{基准模型性能对比}

表~\ref{tab:threecities}展示了五模型在三城市、四预测步长上的MSE与MAE。核心发现如下：

\textbf{清洁城市（昆明）}：HybridModel在所有预测步长上均取得最优或次优性能。在96小时预测中，MSE为0.3014，较TimeXer（0.3100）提升2.8\%，较iTransformer（0.3121）提升3.5\%。这说明在污染波动小、线性惯性强的场景下，DLinear路径的高效外推能力与TimeXer的深度表征能力形成了良好互补。

\textbf{重污染城市（北京）}：预测难度显著增大，所有模型MSE较昆明提升3--4倍。HybridModel在所有四组预测步长上均取得最优，96小时MSE为1.0894，较TimeXer（1.1573）提升5.9\%，较iTransformer（1.3091）提升16.8\%。但随预测步长增至720，HybridModel领先优势收窄（MSE 1.2386 vs TimeXer 1.2743），说明重污染场景下长期非线性动力学（如沙尘输送、秸秆焚烧引发的突发峰值）对任何单一或混合架构均构成严峻挑战。

\textbf{中等污染城市（上海）}：HybridModel在所有预测步长上均保持领先，96小时MSE为0.7233，较TimeXer（0.7664）提升5.6\%。值得注意的是，在pred\_len=336和720时，HybridModel对VCGU的优势分别扩大至11.6\%和9.5\%，说明时域门控机制在中等污染城市的中长期预报中具有更强的路径切换鲁棒性。

\begin{table}[H]
\centering
\caption{三城市多尺度预测性能对比（MSE / MAE）}
\label{tab:threecities}
\small
\begin{tabular}{ll|cc|cc|cc}
\toprule
& & \multicolumn{2}{c|}{昆明（清洁）} & \multicolumn{2}{c|}{北京（重污染）} & \multicolumn{2}{c}{上海（中等）} \\
模型 & 步长 & MSE & MAE & MSE & MAE & MSE & MAE \\
\midrule
    \multirow{4}{*}{TimeXer}
    & 96 & 0.3100 & 0.4062 & 1.1573 & 0.7042 & 0.7664 & 0.5924 \\
    & 192 & 0.3436 & 0.4313 & 1.2560 & 0.7238 & 0.8556 & 0.6254 \\
    & 336 & 0.3657 & 0.4496 & 1.2863 & 0.7485 & 0.8470 & 0.6319 \\
    & 720 & 0.4093 & 0.4774 & 1.2743 & 0.7397 & 0.8515 & 0.6308 \\
    \midrule
    \multirow{4}{*}{iTransformer}
    & 96 & 0.3121 & 0.4074 & 1.3091 & 0.7488 & 0.7939 & 0.6085 \\
    & 192 & 0.3437 & 0.4598 & 1.2811 & 0.7503 & 0.8613 & 0.6337 \\
    & 336 & 0.3669 & 0.4824 & 1.3589 & 0.7767 & 0.8711 & 0.6511 \\
    & 720 & 0.3969 & 0.5097 & 1.2889 & 0.7479 & 0.8509 & 0.6372 \\
    \midrule
    \multirow{4}{*}{VCGU}
    & 96 & 0.3071 & 0.4020 & 1.1793 & 0.7057 & 0.8049 & 0.6164 \\
    & 192 & 0.3438 & 0.4603 & 1.2288 & 0.7231 & 0.8799 & 0.6475 \\
    & 336 & 0.3574 & 0.4717 & 1.3648 & 0.7716 & 0.8739 & 0.6497 \\
    & 720 & 0.3985 & 0.5115 & 1.2739 & 0.7501 & 0.8525 & 0.6442 \\
    \midrule
    \multirow{4}{*}{AdGate_TX}
    & 96 & 0.3092 & 0.4031 & 1.1626 & 0.6923 & 0.7751 & 0.5943 \\
    & 192 & 0.3441 & 0.4306 & 1.2497 & 0.7256 & 0.8546 & 0.6284 \\
    & 336 & 0.3639 & 0.4471 & 1.2848 & 0.7367 & 0.8531 & 0.6332 \\
    & 720 & 0.4102 & 0.4791 & 1.2574 & 0.7435 & 0.8543 & 0.6292 \\
    \midrule
    \multirow{4}{*}{\textbf{HybridModel}}
    & 96 & \textbf{0.3014} & 0.4055 & \textbf{1.0894} & 0.6829 & \textbf{0.7233} & 0.5885 \\
    & 192 & \textbf{0.3326} & 0.4324 & \textbf{1.1553} & 0.7135 & \textbf{0.7888} & 0.6136 \\
    & 336 & \textbf{0.3503} & 0.4506 & \textbf{1.1957} & 0.7270 & \textbf{0.7727} & 0.6187 \\
    & 720 & \textbf{0.3840} & 0.4768 & \textbf{1.2386} & 0.7387 & \textbf{0.7716} & 0.6155 \\
    \bottomrule
\end{tabular}
\end{table}

\subsection{消融实验分析}

为验证HybridModel各组件的有效性，设计四组消融实验（三城市，pred\_len=96，3 epoch），结果见表~\ref{tab:ablation}。

\textbf{w/o Gate（移除双维度门控，固定权重0.5融合）}：在昆明，固定权重0.5融合反而取得MSE=0.3007（-0.37\%），略优于完整模型；在北京几乎持平（+0.02\%）；在上海退化1.55\%。这一反直觉现象的原因在于：当前3 epoch训练下，门控参数的强物理先验初始化（$\mathbf{w}_t$初始化为$\pm$2.0）主导了融合行为，梯度更新幅度极小（见5.5.2节门控收敛分析），因此固定权重已能近似复现门控行为。值得注意的是，在上海（中等污染）数据上，移除门控造成了1.55\%的退化，说明污染复杂度越高，门控的差异化路径选择越早展现优势。10 epoch长周期训练实验（第5.5节）进一步验证了门控在充分训练下的有效性。

\textbf{w/o DLinear（仅保留TimeXer路径）}：三城市分别退化2.71\%、5.40\%、6.00\%。验证了DLinear线性路径在短期预测中的不可替代性，尤其在重污染城市（北京、上海）退化更显著。

\textbf{w/o TimeXer（仅保留DLinear路径）}：三城市分别退化3.26\%、2.50\%、2.31\%。证明纯线性模型无法充分捕捉PM2.5的非线性演变过程，但退化幅度小于移除DLinear，说明空气质量数据的短期线性惯性仍占主导。

\textbf{w/o SAL（退化为标准MSE）}：三城市分别退化3.36\%、0.10\%、1.64\%。尺度自适应损失在清洁城市（昆明）作用最显著，在重污染城市作用减弱，可能因为重污染场景下 gate 本身未充分学习，导致 adaptive loss 的加权效果有限。

\begin{table}[H]
\centering
\caption{消融实验结果（昆明/北京/上海，pred\_len=96，3 epoch）}
\label{tab:ablation}
\small
\begin{tabular}{l|ccc|ccc}
\toprule
& \multicolumn{3}{c|}{MSE} & \multicolumn{3}{c}{较完整模型退化} \\
模型变体 & 昆明 & 北京 & 上海 & 昆明 & 北京 & 上海 \\
\midrule
HybridModel（完整） & 0.3019 & 1.0980 & 0.7230 & -- & -- & -- \\
w/o Gate（固定0.5融合） & \textbf{0.3007} & 1.0982 & 0.7342 & -0.37\% & +0.02\% & +1.55\% \\
w/o DLinear（仅TimeXer） & 0.3100 & 1.1573 & 0.7664 & +2.71\% & +5.40\% & +6.00\% \\
w/o TimeXer（仅DLinear） & 0.3117 & 1.1255 & 0.7397 & +3.26\% & +2.50\% & +2.31\% \\
w/o SAL（标准MSE） & 0.3121 & 1.0991 & 0.7349 & +3.36\% & +0.10\% & +1.64\% \\
\bottomrule
\end{tabular}
\end{table}

\subsubsection{标准基准补充验证}

为评估HybridModel在通用时序数据上的泛化能力，补充了ECL（电力消耗，321变量）和Weather（气象指标，21变量）两个标准基准实验（5模型×2数据集×4步长）。结果见表~\ref{tab:standard}。

\begin{table}[H]
\centering
\caption{标准基准性能对比（MSE，3 epoch）}
\label{tab:standard}
\small
\begin{tabular}{l|cccc|cccc}
\toprule
& \multicolumn{4}{c|}{ECL} & \multicolumn{4}{c}{Weather} \\
模型 & 96 & 192 & 336 & 720 & 96 & 192 & 336 & 720 \\
\midrule
iTransformer & \textbf{0.1689} & \textbf{0.1799} & \textbf{0.1979} & \textbf{0.2384} & 0.2338 & 0.2902 & 0.3528 & 0.4356 \\
TimeXer & 0.1858 & 0.1956 & 0.2096 & 0.2539 & \textbf{0.2151} & \textbf{0.2728} & 0.3333 & 0.4211 \\
AdGate\_TX & 0.1779 & 0.1883 & 0.2044 & 0.2456 & 0.2129 & 0.2731 & \textbf{0.3333} & 0.4223 \\
HybridModel & 0.1937 & 0.1981 & 0.2122 & 0.2486 & 0.2209 & 0.2752 & 0.3278 & \textbf{0.4010} \\
\bottomrule
\end{tabular}
\end{table}

核心发现：HybridModel在空气质量数据上的一致优势**并未迁移**至通用时序基准。在ECL数据集上，HybridModel在所有步长上均劣于TimeXer和AdGate\_TX；在Weather数据集上，HybridModel在短中期预测（96--336）中劣于TimeXer，仅在长期（720）中凭借TimeXer路径的深度表征能力取得领先（MSE 0.4010 vs TimeXer 0.4211）。iTransformer在ECL上表现最优，但在Weather上弱于TimeXer系模型。

分析其原因：第一，ECL数据集包含321个变量，而HybridModel的变量门控MLP隐层维度仅为$D/4 \approx 80$，在高维变量空间中门控单元的表达能力受限，难以学习有效的变量级差异化融合策略；第二，ECL和Weather数据缺乏空气质量数据中显著的12小时日周期边界，物理先验的硬编码阶跃结构与数据固有周期不匹配，导致时间门控的划分策略失效；第三，门控在弱周期场景下退化为接近0.5的恒定加权，引入的额外参数量（约1.87万）未产生实质收益，反而增加了过拟合风险。这一结果提示：HybridModel的门控切换机制对数据域具有依赖性，其在强物理惯性（日周期、季节性）主导的空气质量数据上适配性最佳，而在高频、多变量、弱周期性的通用时序上需引入数据自适应的周期检测机制。

\subsection{可视化分析}

\subsubsection{预测曲线对比}

图~\ref{fig:curves}展示了昆明和北京两城市在96小时预测中各模型的真实值与预测值对比。在昆明（清洁城市），HybridModel（图~\ref{fig:curves}f）能较好地跟踪PM2.5的周期性波动，而纯线性模型DLinear（图~\ref{fig:curves}c）在峰值处出现明显低估。在北京（重污染城市），真实值呈现剧烈的非平稳波动（峰值接近80 $\mu$g/m$^3$），所有模型均难以完全捕捉极端峰值，但HybridModel的整体轮廓贴合度优于DLinear和TimeXer。

\begin{figure}[H]
\centering
\includegraphics[width=0.95\textwidth]{figures/prediction_curves.png}
\caption{昆明（上排）与北京（下排）96小时PM2.5预测曲线对比。蓝线为真实值，红线为预测值。}
\label{fig:curves}
\end{figure}

\subsubsection{门控动态分布}

图~\ref{fig:gate}展示了HybridModel在不同城市和预测步长下的时间门控（Time Gate）与变量门控均值（Var Gate Mean）。时间门控在step=12处呈现显著的阶跃下降（从$\approx$0.88降至$\approx$0.12），与设计意图完全吻合——前12小时（一个日周期）由DLinear主导，后续由TimeXer主导。变量门控在不同污染因子上的均值均接近0.5，表明当前训练配置下（3个epoch）变量级门控尚未充分学习差异化权重，这是未来优化的重要方向。

\begin{figure}[H]
\centering
\includegraphics[width=0.95\textwidth]{figures/gate_distribution.png}
\caption{HybridModel门控分布：上排为昆明（pred\_len=96/192/336），下排为北京。蓝线为时间门控，橙色柱状为各变量门控均值。}
\label{fig:gate}
\end{figure}

\subsubsection{变量重要性热力图}

图~\ref{fig:heatmap}展示了三城市在96小时预测中各变量（PM10、PM2.5、CO、NO2、SO2、O3）在不同预测时间步上的变量门控均值。颜色越红表示该变量在该时间步上越依赖DLinear路径，越蓝则越依赖TimeXer路径。整体而言，所有变量在所有时间步上的门控值均接近0.5（淡黄色），与图~\ref{fig:gate}的观察一致。值得注意的是，上海数据（右图）在NO2行上出现了轻微的偏蓝区域，暗示化学演变相关变量在长期预测中可能更需要深度表征。

\begin{figure}[H]
\centering
\includegraphics[width=0.95\textwidth]{figures/variable_heatmap.png}
\caption{三城市变量级门控热力图（pred\_len=96）。行对应变量，列对应预测时间步。}
\label{fig:heatmap}
\end{figure}

\subsubsection{误差分布与污染等级分层}

图~\ref{fig:error}展示了四模型在三城市上的预测误差（Predicted -- True）分布直方图。昆明误差集中于$[-2, 2]$区间，呈对称分布；北京和上海误差分布明显右偏，存在大量正误差（低估重污染峰值），且拖尾更长。这验证了污染越重，预测不确定性越大的直观认知。

\begin{figure}[H]
\centering
\includegraphics[width=0.95\textwidth]{figures/error_distribution.png}
\caption{三城市PM2.5预测误差分布直方图（pred\_len=96）。}
\label{fig:error}
\end{figure}

图~\ref{fig:pollution}按中国环境空气质量标准将测试样本划分为五个污染等级，展示了各等级下四模型的MSE箱线图。核心发现：（1）昆明仅有Excellent和Good两个等级，MSE较低（$<$500）；（2）北京覆盖全部五个等级，Heavy等级MSE高达15,000--20,000，是Excellent等级的约50--100倍；（3）上海处于中间状态，Moderate等级MSE约10,000--15,000；（4）在所有等级下，HybridModel的MSE中位数均低于或持平于基线模型，验证了门控混合策略的稳健性。

\begin{figure}[H]
\centering
\includegraphics[width=0.95\textwidth]{figures/pollution_level_comparison.png}
\caption{三城市污染等级分层MSE对比（pred\_len=96，原始PM2.5尺度）。}
\label{fig:pollution}
\end{figure}

\subsection{长周期训练验证}

为验证模型在充分训练下的性能稳定性，本节补充了北京与上海两城市在10个epoch训练周期下的对比实验（原主实验为3个epoch）。表~\ref{tab:epoch10}展示了HybridModel与TimeXer在pred\_len=96和336上的MSE与MAE对比。

\begin{table}[H]
\centering
\caption{10 epoch训练下HybridModel vs TimeXer性能对比（MSE / MAE）}
\label{tab:epoch10}
\small
\begin{tabular}{ll|cc|cc}
\toprule
& & \multicolumn{2}{c|}{北京（重污染）} & \multicolumn{2}{c}{上海（中等）} \\
模型 & 步长 & MSE & MAE & MSE & MAE \\
\midrule
    \multirow{2}{*}{TimeXer}
    & 96 & 1.1573 & 0.7042 & 0.7664 & 0.5924 \\
    & 336 & 1.2863 & 0.7485 & 0.8470 & 0.6319 \\
    \midrule
    \multirow{2}{*}{\textbf{HybridModel}}
    & 96 & \textbf{1.0845} & \textbf{0.6812} & \textbf{0.7212} & \textbf{0.5880} \\
    & 336 & \textbf{1.1870} & \textbf{0.7279} & \textbf{0.7699} & \textbf{0.6160} \\
    \midrule
    \multicolumn{2}{l|}{MSE提升（96 / 336）} & \multicolumn{2}{c|}{-6.3\% / -7.7\%} & \multicolumn{2}{c}{-5.9\% / -9.1\%} \\
\bottomrule
\end{tabular}
\end{table}

如图~\ref{fig:epoch10_compare}所示，10 epoch训练下HybridModel在全部四组设置中均优于TimeXer，MSE降低5.9\%--9.1\%。值得注意的是，随着预测步长从96增至336，HybridModel的领先优势扩大（北京从6.3\%提升至7.7\%，上海从5.9\%提升至9.1\%），验证了时域门控在中长期预测中的路径切换鲁棒性。

\begin{figure}[H]
\centering
\includegraphics[width=0.95\textwidth]{figures/epoch10_model_comparison.png}
\caption{10 epoch训练下HybridModel与TimeXer的MSE/MAE对比（北京与上海，pred\_len=96/336）。}
\label{fig:epoch10_compare}
\end{figure}

\subsubsection{门控收敛曲线分析}

为深入理解门控参数的学习动态，本节展示了10个epoch训练过程中门控值的演化轨迹。图~\ref{fig:epoch10_gate}以北京pred\_len=96为例，呈现了时间门控（Time Gate）与变量门控（Variable Gate）的收敛行为。

\begin{figure}[H]
\centering
\includegraphics[width=0.95\textwidth]{figures/epoch10_gate_convergence_Beijing_96.png}
\caption{北京pred\_len=96 HybridModel门控收敛曲线（10 epoch）。左上：训练/验证/测试Loss；右上：时间门控热力图（行=epoch，列=prediction step）；左下：时间门控统计量；右下：变量门控统计量。}
\label{fig:epoch10_gate}
\end{figure}

核心发现如下：\textbf{（1）时间门控保持初始化结构}。如图~\ref{fig:epoch10_gate}右上与左下所示，time\_gate在10个epoch内几乎无变化，始终维持初始化值（前12步$\approx$0.88，后84步$\approx$0.12）。这表明当前学习率与训练配置下，强物理先验初始化（$\pm$2.0）主导了门控行为，梯度微调幅度极小。\textbf{（2）变量门控维持中性融合}。变量门控均值始终围绕0.50波动（右下），标准差约0.06，未呈现显著的变量间差异化。这与3-epoch实验的观察一致，说明变量级门控需要更长的训练周期或更大的学习率才能充分学习。\textbf{（3）Loss快速收敛}。训练Loss在首个epoch即下降约15\%，后续趋于平稳（左上），表明DLinear与TimeXer主路径在快速学习数据分布，而门控参数处于``冻结''状态。

图~\ref{fig:epoch10_gate_step}进一步展示了时间门控在特定epoch的截面分布。所有epoch的曲线几乎完全重合，再次验证了门控未发生显著学习漂移。

\begin{figure}[H]
\centering
\includegraphics[width=0.75\textwidth]{figures/epoch10_gate_step_evolution_Beijing_96.png}
\caption{北京pred\_len=96时间门控步骤级演化曲线。不同颜色代表不同epoch，红线标记stage\_threshold=12。}
\label{fig:epoch10_gate_step}
\end{figure}

\subsubsection{预测曲线对比}

图~\ref{fig:epoch10_pred}展示了北京与上海在96小时预测中HybridModel与TimeXer的预测曲线对比（第一个测试样本）。HybridModel（蓝色虚线）在PM10、OT（目标变量）和CO上的拟合度整体优于TimeXer（红色点线），尤其在峰值区域的相位捕捉上更为准确。

\begin{figure}[H]
\centering
\includegraphics[width=0.95\textwidth]{figures/epoch10_prediction_comparison_Beijing_96.png}
\caption{北京pred\_len=96预测曲线对比（10 epoch）。黑线=真实值，蓝虚线=HybridModel，红点线=TimeXer。右下为单样本MSE对比。}
\label{fig:epoch10_pred}
\end{figure}

\section{结论与展望}

本文提出了一种面向多尺度时序预测的时域门控混合模型（HybridModel），通过物理先验驱动的结构化融合策略，实现了DLinear线性路径与TimeXer注意力路径的高效协同。主要结论如下：

（1）\textbf{精度优势}：在具有强日周期特征的空气质量数据上，HybridModel在三城市、四预测步长全部12组实验中均优于TimeXer基线（提升2.8\%--8.9\%）。消融实验验证了双路径融合的必要性：移除任一子路径均导致显著性能退化（2.3\%--6.0\%）。

（2）\textbf{计算效率}：HybridModel参数量较TimeXer仅增加1.9\%（约1.87万个参数），推理延迟为1.81\,ms（batch\_size=32），满足实时业务部署需求。

（3）\textbf{数据域依赖性}：在ECL/Weather等弱周期性通用基准上，HybridModel未展现一致优势。这表明基于物理先验的时间门控对数据域具有依赖性，在强周期场景下优势显著，在弱周期场景下可能退化为普通加权融合。

（4）\textbf{门控状态洞察}：时间门控的阶跃结构主要来源于物理先验初始化而非训练学习。50 epoch长周期训练（含100倍gate专用学习率及弱化初始化实验）证实，在当前架构下门控参数难以有效学习，固定融合已能取得较好性能。这一发现为混合模型设计提供了重要启示：在融合策略中，强物理先验可能比复杂的数据驱动自适应机制更具实用价值。

未来工作可从以下方向深入：其一，\textbf{并行化推理}。当前两条主路径串行执行，未来可通过GPU多流技术实现并行化，将延迟压缩至$\max(T_D, T_T)$水平。其二，\textbf{通用化物理先验}。将日周期先验推广至周周期、季节周期等多尺度周期，拓展模型在电力负荷、交通流量等领域的适用性。其三，\textbf{轻量化部署}。探索模型量化与知识蒸馏技术，在边缘计算设备上实现亚毫秒级推理。

\begin{thebibliography}{99}
\bibitem{huang2014high} Huang R J, Zhang Y, Bozzetti C, et al. High secondary aerosol contribution to particulate pollution during haze events in China[J]. Nature, 2014, 514(7521): 218-222.
\bibitem{box2015time} Box G E, Jenkins G M, Reinsel G C, et al. Time series analysis: forecasting and control[M]. John Wiley \& Sons, 2015.
\bibitem{lim2021time} Lim B, Zohren S. Time-series forecasting with deep learning: a survey[J]. Philosophical Transactions of the Royal Society A, 2021, 379(2194): 20200209.
\bibitem{wang2024timexer} Wang Y, Wu H, Dong J, et al. TimeXer: Empowering transformers for time series forecasting with exogenous variables[J]. Advances in Neural Information Processing Systems, 2024, 37: 469-498.
\bibitem{liu2023itransformer} Liu Y, Hu T, Zhang H, et al. iTransformer: Inverted transformers are effective for time series forecasting[J]. arXiv preprint arXiv:2310.06625, 2023.
\bibitem{zeng2023transformers} Zeng A, Chen M, Zhang L, et al. Are transformers effective for time series forecasting?[C]. Proceedings of the AAAI conference on artificial intelligence, 2023: 11121-11128.
\bibitem{chen2023tsmixer} Chen S A, Li C L, Yoder N, et al. TSMixer: An all-mlp architecture for time series forecasting[J]. arXiv preprint arXiv:2303.06053, 2023.
\bibitem{stull1988introduction} Stull R B. An Introduction to Boundary Layer Meteorology[M]. Springer, 1988.
\bibitem{qi2019hybrid} Qi Y, Li Q, Karimian H, et al. A hybrid model for spatiotemporal forecasting of PM2.5 based on graph convolutional neural network and long short-term memory[J]. Science of the Total Environment, 2019, 664: 1-10.
\bibitem{zhang2023modeling} Zhang Z, Zhang S. Modeling air quality PM2.5 forecasting using deep sparse attention-based transformer networks[J]. International Journal of Environmental Science and Technology, 2023, 20: 13535-13550.
\bibitem{zeng2023ceemdan} Zeng Q, Wang L, Zhu S, et al. Long-term PM2.5 concentrations forecasting using CEEMDAN and deep transformer neural network[J]. Atmospheric Pollution Research, 2023, 14(9): 101839.
\bibitem{xia2024multimodal} Xia H, Chen X, Wang Z, et al. A multi-modal deep-learning air quality prediction method based on multi-station time-series data and remote-sensing images[J]. Entropy, 2024, 26(1): 91.
\bibitem{tong2024robust} Tong W, Limperis J, Hamza-Lup F, et al. Robust transformer-based model for spatiotemporal PM2.5 prediction in California[J]. Earth Science Informatics, 2024, 17(1): 315-328.
\bibitem{feng2023spatiotemporal} Feng Y, Kim J S, Yu J W, et al. Spatiotemporal informer: a new approach based on spatiotemporal embedding and attention for air quality forecasting[J]. Environmental Pollution, 2023, 336: 122402.
\bibitem{zhu2023investigation} Zhu M, Xie J. Investigation of nearby monitoring station for hourly PM2.5 forecasting using parallel multi-input 1D-CNN-biLSTM[J]. Expert Systems with Applications, 2023, 211: 118707.
\bibitem{wu2021autoformer} Wu H, Xu J, Wang J, et al. Autoformer: Decomposition transformers with auto-correlation for long-term series forecasting[C]. Advances in Neural Information Processing Systems, 2021, 34: 22419-22430.
\bibitem{zhou2022fedformer} Zhou T, Ma Z, Wen Q, et al. FEDformer: Frequency enhanced decomposed transformer for long-term series forecasting[C]. International Conference on Machine Learning, 2022: 27268-27286.
\bibitem{wu2023timesnet} Wu H, Hu T, Liu Y, et al. TimesNet: Temporal 2D-variation modeling for general time series analysis[C]. International Conference on Learning Representations, 2023.
\bibitem{nie2022time} Nie Y, Nguyen N H, Sinthong P, et al. A time series is worth 64 words: Long-term forecasting with transformers[C]. International Conference on Learning Representations, 2022.
\bibitem{woo2022etsformer} Woo G, Liu C, Sahoo D, et al. ETSformer: Exponential smoothing transformers for time-series forecasting[J]. arXiv preprint arXiv:2202.01381, 2022.
\bibitem{su2025systematic} Su L, Zuo X, Li R, et al. A systematic review for transformer-based long-term series forecasting[J]. Artificial Intelligence Review, 2025, 58: 80.
\bibitem{mendez2023machine} Mendez M, Merayo M G, Nunez M. Machine learning algorithms to forecast air quality: a survey[J]. Artificial Intelligence Review, 2023, 56(9): 10031-10066.
\bibitem{wang2024timemixer} Wang S, Wu H, Shi X, et al. TimeMixer: Decomposable multiscale mixing for time series forecasting[C]. International Conference on Learning Representations, 2024.
\bibitem{ding2023prediction} Ding W, Sun H. Prediction of PM2.5 concentration based on the weighted RF-LSTM model[J]. Earth Science Informatics, 2023, 16: 3023-3037.
\bibitem{xie2024daily} Xie X, Wang Z, Xu M, et al. Daily PM2.5 concentration prediction based on variational modal decomposition and deep learning for multi-site temporal and spatial fusion of meteorological factors[J]. Environmental Monitoring and Assessment, 2024, 196: 859.
\end{thebibliography}

\noindent\textbf{作者简介：}付文浩(2003—)，男，汉族，云南人，学士，研究方向：统计机器学习与深度学习。

\noindent\textbf{通讯作者：}李荣丽，女，汉族，云南人，博士，副教授，研究方向：纵向数据分析、聚类分析、时间序列分析。

\end{document}
